% Round 18, higher-centers branch. Candidate until independently reviewed.
\section{Complete-intersection ambients and new fourfolds}
\label{r18ca:section}

The preceding product-smoothing calculation uses the projective
ambient space in two distinct ways: its local smooth coordinates,
and the complete-intersection description of the closed strata.
Both properties persist in a smooth complete-intersection ambient
variety with good reduction.  The resulting generic fibers can be
smooth complete intersections that are not isomorphic to any
hypersurface.  This is a change in the covered generic-fiber class,
not a replacement of the model of an already covered hypersurface.

\subsection{Complete native hypotheses}

Fix the original finite valued extension $K/\mathbf Q_p$, its
valuation ring $R$, a uniformizer $\pi$, and finite residue field
$k$.  Let
\[
 \mathcal Y\subset\mathbf P_R^{n+c+1},\qquad n\ge4,\quad c\ge0,
 \tag{CA18.1}\label{r18ca:ambient}
\]
be a smooth projective relative complete intersection of relative
dimension $n+1$ and multidegree $(b_1,\ldots,b_c)$, with every
$b_j\ge1$.  Its geometric fibers are assumed integral.
Write $h=c_1(\mathcal O_{\mathcal Y}(1))$ and
$\delta=\prod_j b_j$, with $\delta=1$ if $c=0$.
Take $s\ge3$ actual homogeneous forms $f_i$ on the displayed
projective space, of degrees $d_i\ge1$, and use the same
symbols for their restrictions to $\mathcal Y$.
Set $D=\sum_i d_i$.

On $Y=\mathcal Y_k$, require that the divisors
$H_i=V(\bar f_i)$ form a geometric simple-normal-crossings
divisor.  More explicitly, every intersection
$H_I=\bigcap_{i\in I}H_i$ for $1\le |I|\le n+1$ is a
smooth complete intersection in the displayed projective space
of its expected dimension $n+1-|I|$, and every intersection of
$n+2$ factors is empty.  The positive-dimensional intersections
are geometrically integral; a zero-dimensional intersection is
a native finite etale scheme, with every one of its geometric
points retained.  All these are the actual complete intersections
of the ambient defining forms and the indicated factor forms.

Choose an actual homogeneous form $Q$ of degree $D$, restricted
to $\mathcal Y$, whose reduction cuts each positive-dimensional $H_I$ transversely
and avoids every zero-dimensional $H_I$.  Assume
\[
       X=V\!\left(f_1\cdots f_s+\pi Q\right)\subset\mathcal Y_K
 \tag{CA18.2}\label{r18ca:generic}
\]
is smooth and geometrically integral.
It is then a smooth complete intersection of multidegree
$(b_1,\ldots,b_c,D)$ in $\mathbf P_K^{n+c+1}$.
All assumptions and degrees are native.  No alteration or
extension is implicit in them.

\subsection{The actual ordered model and its Chow data}

\begin{lemma}[The same model construction in the smooth ambient variety]
\label{r18ca:model}
Blow up $\mathcal Y$ successively along the strict transforms
of the special divisors $H_1,\ldots,H_{s-1}$, using their
native ambient ideals, and take the strict transform of the
family \eqref{r18ca:generic}.  The resulting $\mathcal X/R$
is regular, projective and reduced strictly semistable, with
components $D_1,\ldots,D_s$ and no extra vertical components.
For each nonempty $I$, writing $j_0=\max I$, its stratum
$D_I$ is the reverse-order iterated blowup of $H_I$ in
the strict transforms of
\[
        Z_{I,j}=H_I\cap H_j\cap V(\bar Q),
                 \qquad j=s,s-1,\ldots,j_0+1.
 \tag{CA18.3}\label{r18ca:centers}
\]
Empty centers are omitted.  Every occurring center has
codimension two in the current stratum, is smooth, and
is isomorphic to its original complete intersection.
The ambient hyperplane class restricts throughout this
identification to its original hyperplane class.
\end{lemma}

\begin{proof}
At a point of $Y$, the factor equations through it are
independent local parameters by the geometric SNC hypothesis.
If $\bar Q$ vanishes there, transversality on the common
stratum supplies one more independent parameter.  After
trivializing the relevant line bundles and absorbing the
unit factors, the actual etale local equation is precisely
\[
                         x_1\cdots x_t=\pi y.
 \tag{CA18.4}\label{r18ca:chart}
\]
If $\bar Q$ is nonzero, its local equation is a unit and
the family is already an ordinary strict crossing.
The ambient variety is smooth over $R$, so these are
coordinates in an actual smooth ambient local chart,
with the remaining parameters carried along.

The ordered ambient blowup ideals are the same $(\pi,x_i)$
ideals as in HP18/NF18.  Their saturated chart calculation
therefore applies without changing any equation: in the
$\pi$-chart substitute $x_i=\pi t_i$ and cancel the common
Cartier factor; in the $x_i$-chart substitute $\pi=x_i t_i$
and make the same cancellation.  Iterating the checked charts
gives the regular reduced crossings and exactly the indexed
old components.  There is no use of a global affine coordinate
system on projective space in this argument.
All centers are native closed subschemes of the smooth
ambient model, so the blowups and the strict transform are
projective.  The flatness and no-extra-component conclusion
are the same local saturation conclusions as in NF18.
The assumed integral generic fiber makes the strict transform
integral.

On $H_I$, the chart transition maps identify the induced
ideals as $(\bar f_j,\bar Q)$, with reverse index order
as in \eqref{r18ca:centers}.  Multiplication of either
local generator by a transition unit does not change
its ideal, so these identifications glue natively.
The original center is a projective smooth complete
intersection of dimension $n-|I|-1$.
The trace of a preceding center on a later center is
the Cartier divisor given by that additional factor.
For a positive-dimensional center the transversality
and expected-dimension hypotheses make it a regular
Cartier divisor whenever nonempty.  A zero-dimensional
center is disjoint from the other centers, since an
intersection of the two centers would place $Q$ on a
zero-dimensional original factor stratum.  Thus the
successive strict center is its blowup in a Cartier
ideal, or in an empty ideal locus, and remains isomorphic
to it.  The local coordinates \eqref{r18ca:chart} give its
regular codimension-two immersion at every stage.
All these maps lie over $Y$, proving the assertion about $h$.
\end{proof}

For an edge $E_{ij}$ with $i<j$, its starting variety
$H_i\cap H_j$ is a smooth complete intersection of dimension
$n-1\ge3$.  The centers $Z_{ij,l}$, $l>j$, are smooth
complete intersections of dimension $n-3\ge1$ and are
geometrically integral.  Weak Lefschetz and the rational
degree computation in NF18 give $N^1(H_i\cap H_j)=\mathbf Qh$.
Consequently the numerical divisor basis is
\[
 N^1(E_{ij})_{\mathbf Q}
    =\mathbf Qh_{ij}\oplus\bigoplus_{l>j}\mathbf Q e_{ij,l},
 \tag{CA18.5}\label{r18ca:edge-divisors}
\]
where each $e_{ij,l}$ is the pulled-back exceptional divisor
of the indicated actual blowup.  The native connectedness
of each center accounts for exactly one such generator.
This is a numerical statement; it does not assert that an
arbitrary actual Chow group is rank one.

The actual triple Gysin columns are
\[
 \tau_{ijk}
   =(d_kh_{ij}-e_{ij,k},\ -d_jh_{ik},\ d_ih_{jk})
       \quad(i<j<k).
 \tag{CA18.6}\label{r18ca:triple}
\]
Indeed the factor divisor on $H_i\cap H_j$ has actual class
$d_kh$, and its own center subtracts one exceptional divisor.
The earlier and later center computations are the local
two-parameter calculation of NF18: the later center and
that strict divisor lie in disjoint opposite charts at the
step where its own center is blown up.  The other two edges
have only the full hyperplane pullback.  Thus higher strata
introduce no missing exceptional correction.
The exceptional block of \eqref{r18ca:triple} is $-I$,
and literal triple Gysin cancellation gives
$\operatorname{im}\tau\subseteq\ker\sigma_2$.

Choose any relatively ample $\mathcal L$ on $\mathcal X$.
Modulo the triple image, multiplication on the edge
fundamental classes by
$H_M=c_1(\mathcal L)+Mh$ has the same matrix $MI+C$
in the remaining hyperplane bases as in NF18.
It is invertible for a sufficiently large nonnegative
integer $M$, which also gives one globally relatively
ample class.  DG18 then proves
\[
                              U_2=0.
 \tag{CA18.7}\label{r18ca:U2}
\]

\begin{theorem}[Full local G for complete-intersection ambient smoothings]
\label{r18ca:full-G}
Under \eqref{r18ca:ambient}--\eqref{r18ca:generic} and all
the stated transversality hypotheses, every specified
geometrically homologically trivial rational cycle on $X$
satisfies strong G in every codimension.
The form is an ordinary actual delta-form on the complement
of that cycle's precise original support, and the equality
holds for every original compactly supported ordinary actual
test.  The assertion applies afresh to all cycles after
every unramified finite valued extension of $K$.
\end{theorem}

\begin{proof}
The argument separates according to parity only to control
the precise numerical and representative gates.

If $n$ is even, the starting edges of dimension $n-1$ and
all their centers of dimension $n-3$ are odd-dimensional
smooth complete intersections.  The all-degree comparison
proved in EV18 therefore applies to both.  Actual Chow
blowup decomposition and its numerical coefficient extraction
give the same edge bases $h^u,e_lh^{u-1}$ in every interior
degree and give numerical-to-geometric comparison on the
final edges.  Multiplication by the actual ambient $h$ sends
each displayed basis to its identically indexed next basis.
The one ample perturbation $H_M$ has square matrices $MI+C_u$
in all these degrees, and one integer $M$ makes them all
invertible.  The proof of EV18 then gives $U_r=0$ for
$1\le r\le n-2$ and $K_{\mathrm{comp},\ell}^a=\ker\rho_a$
in every required degree.  The initial $U_2$ is
\eqref{r18ca:U2}; every other input is identical to that
proof on actual odd-dimensional complete intersections.

If $n=2k+1$ is odd, the starting components and their
codimension-two centers are odd-dimensional complete
intersections.  Hence their numerical groups, with chosen
actual representatives, are precisely the hyperplane and
exceptional-power bases used in OD18.
The edges and their centers are even-dimensional complete
intersections.  Outside their one middle numerical degree,
the same bases exhaust the groups.  In that middle degree
the displayed hyperplane and exceptional subspace has the
diagonal nondegenerate intersection matrix
\[
       \deg_{H_i\cap H_j}(h^{2k}),\qquad
                    -\deg_{Z_{ij,l}}(h^{2k-2}).
 \tag{CA18.8}\label{r18ca:middle-matrix}
\]
The proof is the actual codimension-two blowup pairing:
cross pairings vanish by $p_*1=0$, and exceptional
self-intersection pushes to minus the pairing on the center.
All displayed degrees remain nonzero after replacing the
projective-space starting variety by $H_i\cap H_j$.

The component representatives restrict to geometric classes
in that subspace, because every restricted exceptional divisor
has degree-one cohomology in the native divisor basis, and
its additional factor is the same global hyperplane power.
Its numerical orthogonal complement lies in the literal
Gysin kernel by adjunction, exactly as in OD18.
The projected $MI+C_u$ matrices are again invertible for
one sufficiently large $M$.  Starting with
\eqref{r18ca:U2}, the weaker OD18 induction, using injectivity
only on the restriction image, gives every required $U_r=0$.
Pairing the chosen compatible representative differences
against the actual generators of \eqref{r18ca:middle-matrix}
proves $K_{\mathrm{comp},\ell}^a=\ker\rho_a$ in the remaining
middle degree.  No comparison for arbitrary primitive
middle cycles has been assumed.

Thus in both parities the two DQ17 gates vanish in all
middle cycle codimensions.  ER17 and VT17 kill the cycle
obstruction, and frozen reconstruction gives the asserted
precise-support ordinary actual form and full test quantifier.
The frozen divisor, one-cycle and zero-cycle sectors give
the endpoint codimensions.
An unramified extension preserves the uniformizer exponent
and all geometric complete-intersection hypotheses.
Every comparison and chosen-representative calculation is
then reapplied over its new native residue field and to
its entire new cycle group.
\end{proof}

\subsection{The top monodromy count, including the ambient degree}

\begin{proposition}[Deep strata in a complete-intersection ambient variety]
\label{r18ca:monodromy}
For the preceding model and every $\ell\ne p$,
\[
 \operatorname{rank}N^n
       =\delta\sum_{|I|=n+1}\prod_{i\in I}d_i
                                 -\binom{s-1}{n}.
 \tag{CA18.9}\label{r18ca:rank}
\]
The empty-sum and binomial conventions are as in NG18.
If the number is positive, $X$ has no potential good
reduction, and the frozen monodromy obstruction forces
every strict semistable smooth projective dominant witness
to have crossing depth at least $n$.
\end{proposition}

\begin{proof}
Every positive-dimensional original stratum is geometrically
connected and remains so after the specified blowups.
Every deepest zero-dimensional stratum is unchanged and
has exactly $\delta\prod_{i\in I}d_i$ geometric points,
by its actual transverse complete-intersection degree.
The curves one level above it are unchanged because $Q$
avoids the zero-dimensional factor strata.
Consequently the weight-zero restriction complex is the
simplex cochain complex through degree $n-1$, with
the coordinate of a top subset $I$ repeated
$\delta\prod_{i\in I}d_i$ times.
The rank of its last differential is $\binom{s-1}{n}$,
giving the cokernel dimension in \eqref{r18ca:rank}.

The opposite extreme Gysin matrix is its transpose under
the native degree pairing on the geometric points and
connected curves.  The positive definite rational pairing
on the geometric point basis makes its kernel complementary
to the restriction image.  The extreme $N^n$ is the
identity on the common point summand, so it induces an
isomorphism from that kernel to the restriction cokernel.
These are exactly the primary weight-spectral-sequence
calculations of HM18/NG18; their proof uses the connected
strata and the individual geometric point degrees, not
that the ambient degree is one.
Purity of the proper smooth strata gives degeneration
at $E_2$, and the weights of $H^n$ lie in $[0,2n]$.
Thus $N^n$ kills weights at most $2n-1$ and has image
in weight zero, making the extreme dimension its actual
rank.  The native residue degree of a closed point is
retained as the number of its geometric point summands.
The final assertions are the frozen NM17 obstruction.
Here crossing depth means the largest number of components
meeting simultaneously minus one.
\end{proof}

\subsection{A fourfold class which cannot be a hypersurface}

\begin{corollary}[A full-G complete intersection of type $(2,5)$]
\label{r18ca:quadric-example}
Take $n=4$, let $\mathcal Y\subset\mathbf P_R^6$ be a
smooth relative quadric of dimension five, and take five
degree-one factors and a degree-five $Q$ satisfying all
the hypotheses above.  Then the generic fiber is a smooth
complete-intersection fourfold of type $(2,5)$ with full
local strong G, and
\[
                         \operatorname{rank}N^4=1.
 \tag{CA18.10}\label{r18ca:quadric-rank}
\]
Such data exist at suitable finite places.  No such
geometric generic fiber is isomorphic to any smooth
projective hypersurface fourfold, over any field extension.
\end{corollary}

\begin{proof}
The complete local Green assertion is
Theorem~\ref{r18ca:full-G}.  Here $\delta=2$, $s=5$,
and every factor degree is one, so
\eqref{r18ca:rank} gives $2-\binom44=1$.
In particular this is not a potentially good-reduction
example presented by a more complicated model.

For nonemptiness, choose a smooth quadric over $\mathbf Q$
and five sufficiently general rational hyperplanes.
Their intersections are smooth of the expected dimensions;
the final one consists of two geometric points.
Choose a rational quintic cutting every positive-dimensional
intersection and the ambient quadric smoothly, and avoiding
those final points.  The required choices are nonempty
Zariski open conditions by the characteristic-zero Bertini
theorem for these finitely many very ample systems.
Rational points are dense in the parameter affine spaces,
so the choices can be made over $\mathbf Q$.
Spread the finitely many equations and smoothness conditions
over $\mathbf Z[1/S]$ for a finite set of primes $S$.
The pencil $f_1\cdots f_5+tQ$ has a smooth member at
$t=\infty$, since $Q$ was chosen smooth on the quadric.
Its singular members therefore form a proper closed subset
of the parameter line, and hence a finite set of parameters.
For every sufficiently large prime $p\notin S$ avoiding
that finite set of rational parameters, put $K=\mathbf Q_p$,
$\pi=p$ and $t=p$.  The displayed native special-fiber
conditions and generic smoothness hold.  A smooth positive-
dimensional complete intersection is geometrically connected
and therefore geometrically integral.  This constructs data
of exactly the kind asserted.

We prove the geometric distinction by intrinsic Chern numbers.
Adjunction on a smooth type-$(2,5)$ fourfold in $\mathbf P^6$
gives $K_X=\mathcal O_X(2+5-7)=\mathcal O_X$.
If it were isomorphic to a smooth degree-$e$ hypersurface
$T\subset\mathbf P^5$ over an algebraically closed extension,
adjunction would give $\mathcal O_T(e-6)\simeq\mathcal O_T$.
Intersecting its first Chern class with the third power
of the ample hyperplane class gives $(e-6)e=0$.
Since $e>0$, necessarily $e=6$.

The tangent and normal exact sequences give
\[
 \begin{aligned}
 c(T_X)&=\frac{(1+h)^7}{(1+2h)(1+5h)},&
                c_4(T_X)&=219h^4,&\deg_Xh^4&=10,\\
 c(T_T)&=\frac{(1+h)^6}{1+6h},&
                c_4(T_T)&=435h^4,&\deg_Th^4&=6.
 \end{aligned}
 \tag{CA18.11}\label{r18ca:chern}
\]
For an explicit arithmetic check, the first denominator
inverse through degree four is
$1-7h+39h^2-203h^3+1031h^4$, so its coefficient after
multiplication by $(1+h)^7$ is
$1031-1421+819-245+35=219$.
The sextic coefficient is
$1296-1296+540-120+15=435$.
Thus the intrinsic top Chern degrees, equivalently the
geometric Euler characteristics, are
\[
                   \chi(X)=2190,\qquad \chi(T)=2610.
 \tag{CA18.12}\label{r18ca:euler}
\]
An isomorphism identifies tangent bundles and preserves
their top Chern degrees, contradicting these unequal
values.  The calculation is geometric and is unchanged
by any valued or unvalued field extension.  It proves
the asserted non-hypersurface distinction independently
of any polarization chosen on the hypothetical isomorphism.
\end{proof}

\paragraph{Dependencies and remaining scope.}
The actual model charts and stratum identifications are the
HP18/NF18 charts applied in the stated smooth ambient local
coordinates.  The numerical and representative proofs are
DG18, EV18 and OD18 on the explicitly checked complete
intersections, with the frozen CL17/DQ17/ER17/VT17 interfaces.
The monodromy proof is HM18/NG18 with the native ambient
degree retained.  The Chern calculation comes directly
from the Euler and normal sequences; the classical
complete-intersection Euler formula is recorded, for
example, in J.~Ewing and S.~Moolgavkar,
\emph{Euler characteristics of complete intersections},
Proc.\ Amer.\ Math.\ Soc.\ \textbf{56} (1976), 390--391,
Theorem~1 (the Hirzebruch formula).
This primary paper and the published Macaulay2
\emph{Book3264Examples}, Section 5.4.3, were checked for
the standard Chern/Euler convention.

The theorem covers the specified smoothings, not every
smooth complete intersection.  It does not supply a
number-field full-height assertion without checking all
other finite places, and it does not by itself treat
new cycles after a ramified extension.
In particular it does not infer a good-ambient witness
for an arbitrary variety from the existence of the
good-reduction ambient variety in this particular
explicit construction.  General G remains open.
There is no historical-priority claim.
